\chapter{Modello FSP e proprietà di correttezza}
\label{chap:fsp}
Il modello della versione multithread è nel file \texttt{nbody.lts}.
Le proprietà sono scritte in FLTL e verificate con LTSA.

% ======================================================
%
% 3.1 - FSP MODEL
%
% ======================================================
\section{Modello FSP}
Il modello rappresenta solo la sincronizzazione tra i thread, non i calcoli: il lavoro sul corpo \texttt{b} è descritto dalle azioni \texttt{compute\_velocity[b]} e \texttt{update\_position[b]}.
I processi e le barriere che condividono sono riportati in \cref{fig:struttura_fsp}.

\begin{figure}[H]
  \centering
  \begin{tikzpicture}[
      proc/.style={draw, rounded corners=3pt, fill=gray!8, minimum width=2.6cm,
        minimum height=0.8cm, font=\ttfamily\small},
      bar/.style={draw, rounded corners=3pt, fill=blue!10, minimum width=1.6cm,
        minimum height=0.7cm, font=\ttfamily\small},
      link/.style={thick, gray!70!black}]
    \node[proc] (master) at (0.4, 0) {MASTER};
    \node[bar] (start) at (5, 1.1) {start};
    \node[bar] (phase) at (5, 0) {phase};
    \node[bar] (end) at (5, -1.1) {end};
    \node[draw, dashed, rounded corners=4pt, fit=(start)(end), inner sep=6pt,
      label={[font=\ttfamily\small]below:BARRIERS}] (bars) {};
    \node[proc] (w1) at (10.8, 0.6) {WORKER(1)};
    \node[proc] (w2) at (10.8, -0.6) {WORKER(2)};
    \node[draw, dashed, rounded corners=4pt, fit=(w1)(w2), inner sep=6pt,
      label={[font=\ttfamily\small]below:WORKER\_POOL}] (pool) {};
    \draw[link] (master) -- node[font=\ttfamily\footnotesize, above] {start, end} (bars);
    \draw[link] (pool) -- node[font=\ttfamily\footnotesize, above, align=center] {start, phase,\\end} (bars);
  \end{tikzpicture}
  \caption{Struttura del sistema \texttt{SIM} con due worker: sulle linee le barriere condivise.}
  \label{fig:struttura_fsp}
\end{figure}

Il sistema verificato è la loro composizione:
\begin{lstlisting}[style=fsp]
||SIM = (MASTER || WORKER_POOL || BARRIERS).
\end{lstlisting}

% ======================================================
% 3.1.1 - BARRIERS
% ======================================================
\subsection{\texttt{BARRIERS}: Barriere}
La barriera conta i thread arrivati con \texttt{arrive}; quando ci sono tutti, \texttt{release} li sblocca insieme e il contatore torna a zero:
\begin{lstlisting}[style=fsp]
BARRIER(P=1, F=0) = COUNT[0],
COUNT[c:0..P] =
    (when (c < P)  arrive[t:F..N_WORKERS] -> COUNT[c+1]
    |when (c == P) release -> COUNT[0]
    |when (c == 0) shutdown -> END).
\end{lstlisting}

\begin{figure}[H]
  \centering
  \begin{tikzpicture}[
      st/.style={circle, draw, fill=blue!10, minimum size=0.9cm, font=\small},
      tr/.style={-{Stealth}, thick},
      lbl/.style={font=\ttfamily\footnotesize}]
    \node[st, double] (c0) at (0, 0) {0};
    \node[st] (c1) at (2.4, 0) {1};
    \node[st] (c2) at (4.8, 0) {2};
    \node[st] (c3) at (7.2, 0) {3};
    \node[st, fill=gray!15] (e) at (0, -2) {\scriptsize END};
    \draw[tr] (c0) -- node[lbl, above] {arrive} (c1);
    \draw[tr] (c1) -- node[lbl, above] {arrive} (c2);
    \draw[tr] (c2) -- node[lbl, above] {arrive} (c3);
    \draw[tr] (c3) to[bend left=35] node[lbl, below] {release} (c0);
    \draw[tr] (c0) -- node[lbl, left] {shutdown} (e);
  \end{tikzpicture}
  \caption{Una barriera per tre thread: il numero nello stato è il contatore \texttt{c}.}
  \label{fig:barriera}
\end{figure}

Le barriere sono \texttt{start}, \texttt{phase} ed \texttt{end}, come nel \cref{sec:design_thread}.
Su \texttt{start} il master può anche fare \texttt{shutdown}, che termina tutti i thread e tutte le barriere.

% ======================================================
% 3.1.2 - WORKERS
% ======================================================
\subsection{\texttt{WORKER\_POOL}: Worker}
Il worker \texttt{I} gestisce un blocco di corpi.
Prima calcola le velocità, poi aspetta su \texttt{phase}, poi aggiorna le posizioni e infine aspetta su \texttt{end}:
\begin{lstlisting}[style=fsp]
WORKER(I=1) =
    (start.arrive[I] ->
        (start.release -> VELOCITY[(I-1)*N_BODIES/N_WORKERS]
        |shutdown -> END)),
VELOCITY[b:(I-1)*N_BODIES/N_WORKERS..I*N_BODIES/N_WORKERS] =
    if (b < I*N_BODIES/N_WORKERS)
    then (compute_velocity[b] -> VELOCITY[b+1])
    else (phase.arrive[I] -> phase.release ->
          POSITION[(I-1)*N_BODIES/N_WORKERS]),
POSITION[b:(I-1)*N_BODIES/N_WORKERS..I*N_BODIES/N_WORKERS] =
    if (b < I*N_BODIES/N_WORKERS)
    then (update_position[b] -> POSITION[b+1])
    else (end.arrive[I] -> end.release -> WORKER).
\end{lstlisting}
Il range di \texttt{b} è limitato al blocco del worker.
Con \texttt{0..N\_BODIES}, infatti, LTSA dava a ogni worker le azioni di tutti i corpi e il sistema andava in deadlock.

% ======================================================
% 3.1.3 - MASTER
% ======================================================
\subsection{\texttt{MASTER}: Master}
Il master, cioè il thread che chiama \texttt{step()}, esegue \texttt{STEPS} passi e poi termina i worker:
\begin{lstlisting}[style=fsp]
MASTER[s:STEP_IDX] =
    if (s < STEPS)
    then (start.arrive[MASTER_ID] -> start.release ->
          end.arrive[MASTER_ID] -> end.release ->
          next_step -> MASTER[s+1])
    else (shutdown -> TERMINATED),
TERMINATED = (terminated -> TERMINATED).
\end{lstlisting}
Il ciclo su \texttt{terminated} serve solo perché LTSA non scambi la fine della simulazione per un deadlock.

% ======================================================
%
% 3.2 - CORRECTNESS PROPERTIES
%
% ======================================================
\section{Proprietà di correttezza}
Per scrivere le proprietà servono alcuni fluenti, cioè condizioni che diventano vere o false in base alle azioni.
Un passo inizia con \texttt{start.release}, che riporta a falso i fluenti dei corpi:
\begin{lstlisting}[style=fsp, basicstyle=\ttfamily\footnotesize]
fluent VELOCITY_DONE[b:BODIES] = <compute_velocity[b], start.release>
fluent POSITION_DONE[b:BODIES] = <update_position[b], start.release>
fluent STEP_RUNNING = <start.release, end.release>
\end{lstlisting}
\begin{description}
  \item[\texttt{VELOCITY\_DONE[b]}] la velocità del corpo \texttt{b} è già stata calcolata in questo passo;
  \item[\texttt{POSITION\_DONE[b]}] la posizione del corpo \texttt{b} è già stata aggiornata in questo passo;
  \item[\texttt{STEP\_RUNNING}] un passo è in corso.
\end{description}

% ======================================================
% 3.2.1 - SAFETY PROPERTIES
% ======================================================
\subsection{Proprietà di Safety}
Le proprietà di safety dicono che non succede mai niente di sbagliato.
Gli indici \texttt{b} e \texttt{c} variano su tutti i corpi.
\begin{description}
  \item[\texttt{CONSISTENT\_READ}] le velocità si calcolano con le posizioni del passo precedente:
    \[\Box\,(\texttt{compute\_velocity[b]} \Rightarrow \textstyle\bigwedge_{c} \neg\texttt{POSITION\_DONE[c]})\]
  \item[\texttt{PHASE\_SEPARATION}] le posizioni si aggiornano solo dopo tutte le velocità:
    \[\Box\,(\texttt{update\_position[b]} \Rightarrow \textstyle\bigwedge_{c} \texttt{VELOCITY\_DONE[c]})\]
  \item[\texttt{ONCE\_PER\_STEP}] ogni corpo si aggiorna una sola volta per passo:
    \[\Box\,(\texttt{update\_position[b]} \Rightarrow \mathsf{X}\,(\neg\texttt{update\_position[b]}\ \mathsf{W}\ \texttt{start.release}))\]
  \item[\texttt{CONSISTENT\_STATE}] il master legge lo stato solo a passo finito:
    \[\Box\,(\texttt{next\_step} \Rightarrow (\neg\texttt{STEP\_RUNNING} \land \textstyle\bigwedge_{b} \texttt{POSITION\_DONE[b]}))\]
\end{description}

% ======================================================
% 3.2.2 - PROGRESS PROPERTIES
% ======================================================
\subsection{Proprietà di Progress}
Le proprietà di progress dicono che prima o poi succede qualcosa di buono.
Il sistema termina dopo \texttt{STEPS} passi, quindi il controllo di progress standard di LTSA non ha senso: si usano invece queste proprietà di liveness.
\begin{description}
  \item[\texttt{STEP\_COMPLETION}] ogni passo iniziato finisce:
    \[\Box\,(\texttt{start.release} \Rightarrow \Diamond\,\texttt{next\_step})\]
  \item[\texttt{BODY\_PROGRESS}] in ogni passo ogni corpo viene aggiornato:
    \[\Box\,(\texttt{start.release} \Rightarrow \Diamond\,\texttt{update\_position[b]})\]
  \item[\texttt{SIM\_TERMINATION}] la simulazione termina:
    \[\Diamond\,\texttt{terminated}\]
\end{description}
